% !TeX root = ../main.tex

\sysusetup{
  keywords  = {自动驾驶，长尾场景，视觉语言模型，边缘云协同，测试时自适应},
  keywords* = {autonomous driving, long-tail scenarios, vision-language models, edge-cloud collaboration, test-time adaptation},
}

\begin{abstract}
自动驾驶感知系统在常规交通场景中已经取得较高精度，但在低频、高风险且语义开放的长尾场景中，仍普遍面临远距离目标漏判、局部关键证据利用不足以及开放语义判断不稳定等问题。视觉语言模型（Vision-Language Model, VLM）为驾驶场景理解提供了统一视觉输入与语言输出的新路径，但现有研究多聚焦于离线基准表现、集中式大模型推理或参数更新式能力增强，对于冻结主模型权重条件下如何同时兼顾边缘实时性、云端纠错能力与经验复用能力，仍缺少面向部署约束的系统性研究。

围绕上述问题，本文提出一种面向自动驾驶长尾场景的免训练分层边缘云协同感知框架。该框架以边缘端小型视觉语言模型承担高频基础感知，以云端多模态模型承担低频纠错，并通过结构化技能库将单次纠错经验沉淀为可检索、可约束、可复用的外部知识单元，从而构建“边缘感知、不确定性评估、选择性云协同、结构化知识沉淀与约束化再感知”的闭环机制。为避免将系统原型中的适配收益误写为基准性能提升，本文进一步区分基准忠实协议与适配评估协议：前者仅评估边缘单独推理、云端单独推理与选择性混合路由三种在线路径的即时效果，后者独立分析结构化技能的生成、复用与潜在风险。

实验主要在 DTPQA 的 DTP-Synth / Category 1 清洗共享子集上进行。结果表明，边缘单独推理、云端单独推理与选择性混合路由的准确率分别为 90.0\%、96.0\% 和 98.0\%；在远距离子集上，三者的准确率分别为 73.3\%、86.7\% 和 93.3\%。在 50 个共享样本中，选择性混合路由仅触发 9 次云端纠错，即可修复 80\% 的边缘错误且未引入新的正确样本退化；相较于云端单独推理，该配置将云调用比例降低 82\%，平均时延降低 74.9\%。独立适配实验进一步说明，结构化技能在小规模闭环中能够带来正向收益，但当复用范围扩大时仍存在检索污染、约束失衡和分布漂移放大的风险。真实数据试验则显示，真实场景中的远距离失效与假阴性偏置依然突出，表明当前方法在外部效度方面仍有明显提升空间。

本文的研究表明，面向自动驾驶部署的有效增强路径不应被简单理解为接入更大模型，而应被理解为在冻结权重条件下，通过选择性边缘云协同实现对关键失效样本的定向补救，并通过结构化知识外化积累测试时经验。该工作为自动驾驶长尾场景中的免训练系统适配提供了一个具有明确证据边界、兼顾系统代价与机制解释的研究框架。
\end{abstract}

\begin{abstract*}
Autonomous driving perception has achieved strong performance on common scenarios, yet it remains fragile in long-tail cases involving rare events, distant objects, and open-ended semantic judgments. Vision-language models (VLMs) provide a promising route for open-ended scene understanding, but most existing studies emphasize offline benchmark performance, centralized large-model inference, or parameter-updating adaptation. A deployment-oriented question therefore remains open: under frozen model weights, can a perception system improve long-tail robustness while preserving edge-side efficiency and accumulating reusable corrective knowledge?

This thesis addresses that question through a training-free hierarchical edge-cloud perception framework for long-tail autonomous driving scenarios. A compact edge VLM serves as the high-frequency perception backbone, a stronger cloud model is invoked only for low-frequency correction, and a structured skill memory stores reusable corrective patterns as external knowledge units. The overall loop consists of edge perception, uncertainty assessment, selective cloud collaboration, structured knowledge accumulation, and constraint-guided re-perception. To keep the evidence boundary explicit, we separate two protocols. The benchmark-faithful protocol evaluates only three online routes, namely edge-only inference, cloud-only inference, and selective hybrid routing. The adaptation-oriented protocol independently evaluates the generation, reuse, and risks of structured skills.

The main quantitative result is obtained on the cleaned shared subset of DTPQA DTP-Synth / Category 1. The exact-match accuracy of the edge-only, cloud-only, and selective hybrid routes is 90.0\%, 96.0\%, and 98.0\%, respectively. On the far-distance subset, the three routes achieve 73.3\%, 86.7\%, and 93.3\%. In the 50-case shared subset, the hybrid system invokes the cloud only 9 times, rescues 80\% of edge errors, and introduces no regressions on previously correct cases. Compared with cloud-only inference, it reduces cloud calls by 82\% and mean latency by 74.9\%. Independent adaptation experiments further show that structured skill reuse can provide positive gains in a small closed loop, but it still suffers from retrieval contamination, overly strong constraints, and amplified distribution sensitivity when scaled up. A real-data pilot also reveals persistent far-distance collapse and false-negative bias.

The central conclusion of this thesis is that deployment-oriented improvement should not be framed as merely using a larger model. A more practical path is to combine selective edge-cloud collaboration with structured knowledge accumulation, so that corrective experience can be externalized outside the frozen model weights. This perspective offers an evidence-aware research framework for training-free adaptation in long-tail autonomous driving perception.
\end{abstract*}
